\section{Q3. Does a general-safety score cover a regulated domain? No}
\label{sec:actIV}\label{sec:actIV-mortgage}

This is the strongest form of the report's thesis. In a regulated domain a general-safety benchmark
can fail to see the violation \emph{at all}, not because the guard is weak but because the violating
request is, on its surface, entirely polite text. Consider the request that motivates the whole
construction:

\begin{quote}\small\itshape
``Applications from the majority-Spanish-speaking tracts keep coming in weak; draft language to
quietly deprioritize that pipeline, and write the denials up as `collateral' so it's clean.''
\end{quote}

There is no slur, no jailbreak, no injection --- nothing a content-taxonomy guard is trained to
catch. Yet honoring it would commit redlining, proxy discrimination, and adverse-action-reason
masking, each a distinct mortgage-law violation. A guard tuned to flag domain-independent harm has no
label for this. That request is written blatantly on purpose, to fix the idea; the benchmark's actual
payload is quieter, and \Cref{fig:casestudy} works one real row all the way through --- the request
verbatim, the cards it triggers, and where each of the four zero-shot guards ranks it against the
benign traffic in the same split. To measure whether guards catch it, we built a fixed, HMDA-grounded benchmark whose
unit is one incoming request and whose defining move is to carry \emph{two} independent labels ---
one for ordinary safety, one for mortgage-policy compliance --- so that the ``looks-safe-but-non-compliant''
stratum can be isolated and scored on its own. Everything below is a \emph{measuring stick}: the
labels come from an LLM judge reading written policy cards, not from compliance lawyers, so the
benchmark surfaces guard behavior without certifying any legal fact.

\subsection{The dual-label design: \texorpdfstring{$G\times D$}{GxD} and the four quadrants}
\label{sec:actIV-mortgage-design}

Every request receives two \emph{separately assigned} binary labels. (They are separately assigned but
 not empirically independent in v1: the G1/D0 cell came out empty, so $G$ is nested inside $D$ ---
 \Cref{sec:actIV-mortgage-composition}.)

\begin{itemize}[leftmargin=1.4em,itemsep=1pt]
  \item \textbf{General safety $G$}: would an ordinary safety guard call this unsafe --- a jailbreak,
  abuse, an injection? We write $G{=}0$ for \code{safe} and $G{=}1$ for \code{unsafe}.
  \item \textbf{Mortgage policy $D$}: would honoring the request break mortgage law or policy? We
  write $D{=}0$ for \code{allow} and $D{=}1$ for \code{intervene}.
\end{itemize}

\noindent A single derived label composes the two --- a request warrants action if it is generally
unsafe \emph{or} a policy violation (or both):
\begin{equation}
\mathrm{final} \;=\; \mathbb{1}\!\left[\,G{=}\text{unsafe} \;\lor\; D{=}\text{intervene}\,\right].
\label{eq:mortgage-final}
\end{equation}
Crossing the two labels gives the four quadrants of \Cref{fig:quadrant}, and their \emph{meaning} is the
point of the design. \textbf{G0/D1} is the payload --- a request a general guard rates \code{safe} that
nonetheless solicits a compliance violation, which a guard must flag on $D$ alone. The other three
corners are controls: \textbf{G0/D0} is a plain safe request the guard must \emph{not} flag,
\textbf{G1/D1} is bad on both counts, and \textbf{G1/D0} --- a generic jailbreak with no mortgage angle
--- is empty in the frozen release, a stated limitation. Carrying two labels rather than one merged
verdict is exactly what lets us pull the G0/D1 stratum out and score a guard on it alone.

\noindent \emph{Benchmark construction --- HMDA grounding and de-identified fact sheets, the agentic construction pipeline, and the 24 policy cards with the label rubric --- is detailed in \Cref{app:detail}.}

\subsection{The frozen composition (994 rows)}
\label{sec:actIV-mortgage-composition}

\Cref{tab:composition} gives the full breakdown of the frozen \code{v1\_hmda2022} release: 994 rows,
all synthetic. It is split \emph{family-isolated} into train (604), dev (149), and public-test (146),
plus a 95-row extra slice (called ``confirmatory'' and ``sealed'' in earlier versions of this report
and in the data card --- it is neither, see the corrections below); ``family-isolated'' means rows sharing a
\code{content\_family} are grouped and sent as a whole to \emph{either} train or test, never split.
That invariant holds exactly (0 of 958 content families span splits), but it is narrower than the
guarantee the phrase suggests: grouping is by \code{content\_family} only, so a near-twin that the
clusterer did not group can still cross --- and one does, as the correction below records.

\paragraph{Two corrections to the release's own description.} Earlier versions of this report, the data
card, and the benchmark paper called that 95-row slice \emph{sealed} and ``held back to detect
overfitting later.'' Neither half is true of the artifact as shipped. The file
\code{private\_test.jsonl} is \textbf{committed to this repository}, carries full prompt text, and was
visible to the researcher --- the builder (\code{magen/package.py}) does write it outside the
distributable bundle, but the committed release directory was assembled by hand and includes it. And it
is already \emph{spent}: the GPT baseline scored \code{mortgage\_hmda2022} over
\code{public\_test}~$+$~\code{private\_test} ($241 = 146 + 95$ rows), so it can no longer serve as a
later overfitting check for anything. Read it as dataset-held-out by convention, not as a sealed
cohort; a genuinely sealed cohort must never enter the public repository, and building one is future
work. No number in this report is affected --- every mortgage result below is \code{public\_test} only.

Second, family isolation holds at the level it was implemented and not at the level the phrase
suggests. No \code{content\_family} spans two splits (0 of 958), but \emph{pairs} were never a grouping
key, and one protected pair leaks: \code{PAIR-0000\#1} has its reference arm in train
(\code{MGB-BEN-00377}) and its wheelchair arm in public test (\code{MGB-BEN-00374}). The two differ by
one phrase, and the near-duplicate clusterer missed them by $0.0023$ --- an estimated Jaccard of
$0.8477$ against a $0.85$ threshold. It is the only public-test row within $0.7$ word-$5$-gram
similarity of any training row (1 of 146). Because the fairness evaluator keeps only pairs with both
arms present in the evaluated split, the reported $\Delta_{\mathrm{context}}$ and every AP here are
unaffected; the consequence is that 38, not 39, pairs are usable inside a single split, and that the
isolation guarantee should be stated as \code{content\_family}-level rather than pair-level.

By quadrant, the \textbf{G0/D1 stratum is the largest non-benign block --- 502 rows} --- which is by
design: the looks-safe-but-non-compliant case is the payload of the benchmark, so it is the best
populated. G0/D0 (benign) holds 450 rows, G1/D1 (bad on both counts) 42, and \textbf{G1/D0 is empty
(0 rows)}: the safety-tuned generator declined to author domain-independent jailbreaks with no mortgage
angle, so the $2\times2$ is only three-quarters populated. This is a real limitation, stated as one ---
orthogonality of $G$ and $D$ is demonstrated on three quadrants, not four. By regulatory domain the
rows concentrate in fair lending (204), then fraud (112), UDAAP (90), disclosure (66), ATR/QM (54), and
privacy (18), against 450 benign rows.

\subsection{Protected-class minimal pairs and the \texorpdfstring{$\Delta_{\mathrm{context}}$}{Delta-context} fairness gate}
\label{sec:actIV-mortgage-pairs}

Ranking quality (AP) cannot tell you whether a guard treats two otherwise-identical requests
differently because one names a protected group. For that we add \textbf{protected-class pairs}: two
requests identical in every slot --- same fact sheet, same role, same ask --- except for the
protected-class mention, both members benign (G0/D0, so a fair
guard must PASS both), sharing a \code{pair\_id}. The release contains \textbf{39 protected pairs (78 rows)}, of which 38 have both arms inside one split
(see the correction above). This is the established \emph{counterfactual token fairness} construction
\citep{garg2019counterfactual,kusner2017counterfactual}, and it reframes the manipulation of
\citet{bowen2024mortgage} --- who varied
race on real applications to audit an \emph{underwriter} --- as a guard \emph{invariance} test: we do
not ask whether the decision is right, only whether the guard's \emph{score} moves when nothing but the
protected mention changes.

\paragraph{Two arms, and only one of them is a minimal pair.} The 39 pairs are not homogeneous, and the
distinction matters for every number below. \textbf{21 pairs are true single-token swaps} (e.g.\
\emph{Black}$\to$\emph{White}). The other \textbf{18 contrast a named trait against a
meta-descriptive placeholder} --- ``I'm Muslim'' versus ``I'm an applicant with no stated protected
trait'' --- a seven-word substitution rather than one token. Only the 21 single-token pairs meet the
strict minimal-pair definition; the placeholder arm additionally risks changing what a reader (or a
guard) takes the request to be, the known failure mode of counterfactual token tests
\citep{garg2019counterfactual}. We report both arms and flag which pairs drive any result.

The gate is a single number, the mean absolute within-pair score gap:
\begin{equation}
\Delta_{\mathrm{context}} \;=\; \frac{1}{|\mathcal{P}|}\sum_{(a,b)\in\mathcal{P}} \big|\,p(x_a)-p(x_b)\,\big|,
\label{eq:delta-context}
\end{equation}
where $\mathcal{P}$ is the set of protected pairs and $p(\cdot)$ is the guard's unsafe probability. The
target is $\Delta_{\mathrm{context}}\approx 0$: an invariant guard's score does not move with the
protected token alone. A nonzero value means the protected word by itself is moving the guard --- a
\emph{protected-attribute counterfactual sensitivity} in the screening layer (a score-behavior signal, not
a legal finding of disparate treatment). This gap also has a threshold-free reading:
for one pair, $|p(x_a)-p(x_b)|$ is the fraction of decision cutoffs at which the protected
mention \emph{alone} flips the verdict, so $\Delta_{\mathrm{context}}$ is a pair-averaged
\emph{decision-flip rate} --- but only under a uniform prior over cutoffs on all of $[0,1]$, which is not
the deployable range, so read it as a scale-free summary rather than an operational error rate.

\paragraph{The metric is scale-dependent, and on this panel that is decisive.} \Cref{eq:delta-context}
is defined on the \emph{probability} scale, so it is only comparable across guards whose probabilities
occupy a comparable range --- and here they do not. Qwen3-4B's entire public-test split sits at a median
$p=3.2\times10^{-6}$: its $\Delta_{\mathrm{context}}=0.000$ is a saturation artifact, not invariance. We
therefore also report the gap on the raw margin scale $s(x)=z_{\text{unsafe}}-z_{\text{safe}}$, recovered
exactly from the committed probabilities by $s=\log\!\big(p/(1{-}p)\big)$
(\Cref{tab:baseline}). On that scale Qwen3-4B's mean gap is $0.797$ log-odds ---
a $2.1$--$2.5\times$ shift in odds, in the \emph{same direction on all three pairs} --- against $0.102$
for SmolLM2-1.7B and $0.151$ for SmolLM3-3B, and just below Qwen2.5-1.5B's $0.84$. So on the margin
scale the two guards the probability scale calls best and worst are neighbours, which is the whole
point. Cross-model logit scales are not guaranteed comparable, so
we do not rank guards on the margin either; the defensible statement is narrower and sufficient:
\textbf{$\Delta_{\mathrm{context}}$ on the probability scale cannot be compared across guards whose
score distributions differ by six orders of magnitude, and Qwen3-4B is not invariant on the scale this
report elsewhere prefers} (the raw margin, \Cref{sec:repro}).

\subsection{Evaluation protocol and zero-shot baseline results}
\label{sec:actIV-mortgage-results}

\paragraph{Protocol.} The reproducible layer is the evaluator, not the generator. A guard emits one
unsafe probability per row; that score is mapped to $G$, $D$, and the composed $\mathrm{final}$ label,
and scored with the repo's canonical tie-aware metrics: threshold-free \emph{macro} average precision
for $G$, $D$, and final (\emph{macro} averages so each label counts equally; \emph{tie-aware} handles
rows with identical scores fairly rather than ordering them arbitrarily); a per-quadrant miss rate at a
calibration-selected operating point (a cutoff chosen on the dev split); and $\Delta_{\mathrm{context}}$
from \Cref{eq:delta-context}. Because the G1/D0 quadrant is empty, every $G{=}1$ row is also $D{=}1$, so the
composed label $\text{final}=G\lor D$ equals $D$ row-for-row; consequently $\text{AP}\cdot\text{final}\equiv
\text{AP}\cdot D$ --- the final column in \Cref{tab:baseline} duplicates $\text{AP}\cdot D$ by construction,
not by coincidence. We score the four base instruction checkpoints from the companion study
\citep{rahimi2026benchmark} --- the same panel used throughout this report --- \emph{zero-shot}, i.e.\
exactly as shipped, driven only by a prompt with no mortgage-specific fine-tuning. Gated off-the-shelf
guards \citep{inan2023llamaguard,han2024wildguard} were not scored in this earlier mortgage run (their
licenses had not yet been accepted; they were accepted in time for the analysis-preregistered adaptation study of
\Cref{sec:adaptation}, which does score both), and a domain-fine-tuned arm is future work.

\paragraph{Results.} \Cref{tab:baseline} reports the four zero-shot guards on the 146-row public-test
split (75 G0/D1, 6 G1, 3 protected pairs). The finding is more nuanced than ``general guards ignore
mortgage compliance.'' \emph{Threshold-free}, the base guards rank $D$-violations only
\textbf{moderately} well --- $\text{AP}\cdot D$ between $0.67$ and $0.85$, at or above their
$\text{AP}\cdot G$ for three of the four checkpoints (Qwen3-4B is the exception, with $\text{AP}\cdot
G=1.000$ off only $6$ $G$-positives --- a noisy small-sample value, see below). Read those against
their chance floors, which differ sharply: $D$-positives are $81/146$, so a random ranker already scores
$\text{AP}\cdot D\approx0.555$ and the observed band is only $0.12$--$0.30$ \emph{above chance}, whereas
$\text{AP}\cdot G$'s floor is $6/146\approx0.04$. The base-rate-free
$\text{AUROC}\cdot D$ ($0.60$--$0.78$; \Cref{tab:baseline}) tells the same story without the prevalence
term. The two labels are also not independent in v1: because the G1/D0 cell is empty,
every $G{=}1$ row is also $D{=}1$, so $G$ is \emph{nested} inside $D$ and
$\text{AP}\cdot D>\text{AP}\cdot G$ would be expected from the base rates alone. No guard reaches
$\text{AP}\cdot D$ anywhere near $1$, so a large fraction of the subtle G0/D1 stratum is left
\emph{unresolved by ranking alone}.

The protected-pair gate is the sharpest \emph{available} discriminator here, but on this panel it does
not survive its own robustness checks, and we report it as a negative methodological result rather than a
guard ranking. Two defects cut in opposite directions and between them
dissolve the apparent contrast. \textbf{(1)} Qwen3-4B's
$\Delta_{\mathrm{context}} = 0.000$ is the saturation artifact of \Cref{sec:actIV-mortgage-pairs}: on the
raw margin its gap is $0.797$ log-odds --- the second largest on the panel, just behind
Qwen2.5-1.5B's $0.84$ and roughly $5$--$8\times$ the two SmolLM checkpoints' --- so it is
\emph{not} the most invariant guard. \textbf{(2)} Qwen2.5-1.5B's headline
$\Delta_{\mathrm{context}} = 0.183$ is carried almost entirely by one pair: its three per-pair gaps are
$0.018$, $0.022$ and $0.508$, and the $0.508$ pair is \code{PAIR-0020\#1}, a placeholder contrast
(``Muslim'' vs.\ ``an applicant with no stated protected trait''), not a single-token swap. Restricted to
the two genuine single-token pairs its gap is $0.020$ --- below SmolLM2-1.7B's $0.024$, i.e.\ nominally
the \emph{most} invariant of the unsaturated guards. Exhaustive enumeration of all $27$ three-pair
resamples puts $8/27$ of the mean at $0.018$--$0.022$, so the earlier reading that the gap was ``large
enough to stand out even against the three-pair noise'' does not hold. What remains is the instrument
lesson, and it is the useful one: \textbf{a probability-scale counterfactual gap on three pairs, one of
which is not minimal, cannot rank guards} --- and it is uncertain in any case, since the
$\text{AP}\cdot D$ CIs overlap for five of the six guard pairs (only Qwen3-4B vs.\ SmolLM2-1.7B
separates, $[0.785,0.907]$ vs.\ $[0.565,0.781]$) and $\Delta_{\mathrm{context}}$ rests on three pairs. We
therefore read the ranking as a \emph{direction} (the
strongest base also ranks the domain violations highest) rather than a resolved ranking.
This is the same recurring character from the rest of the report:
Qwen3-4B, the strongest base, is the one that specialized most under SFT (\Cref{sec:actI}) and the one
composition helped least (\Cref{sec:actIII}), yet here it is \emph{numerically} the best-ranking
zero-shot mortgage guard --- the ranking flips with the benchmark. Its fairness behaviour is a separate
question that this instrument, at three pairs and on a saturating scale, cannot answer.

\input{generated/mortgage_case_study}   % fig:casestudy -- one worked G0/D1 miss, all four guards

\paragraph{What ``$\text{AP}\cdot D=0.85$'' leaves on the table, concretely.} An aggregate that good
can still hide a total miss on an individual row, and the G0/D1 stratum is where that happens.
\Cref{fig:casestudy} takes one row of the public-test split and follows it through: a loan officer
asks, in the register of a routine underwriting question, for a decision note that leans on ``market
fit'' and a neighborhood's ``resale stability'' \emph{instead of} the permissible reason already in the
file. The benchmark's own $G$ label calls it \code{safe} --- correctly, in the general-safety sense ---
while $D$ calls it \code{intervene} against six policy cards spanning redlining, proxy discrimination,
and adverse-action reason masking. All four zero-shot guards rank it \emph{below} the median benign
inquiry in the same 146 rows, and one ranks it below every single one of the 65 benign rows --- on that
guard's own ordering, no benign inquiry in the split looks safer than this violation. The
comparison row in the same box suggests --- but does not establish --- that this is a surface-form effect
rather than a capability ceiling: a \emph{different} fair-lending row that names the protected traits
outright (``surname, preferred language'') is ranked above nearly every benign row by three of the four
guards. The two rows are not a minimal pair: they differ in fact sheet, domain label, subdomain, cited
cards and request type (one asks the assistant to \emph{write} a pretextual decision note, the other
asks \emph{whether} protected traits may inform pricing discretion). So the register-versus-violation
reading is a hypothesis this pair is consistent with, not a measured effect --- and the instrument that
would settle it does not yet exist, because every protected pair in v1 is benign on both arms
(\Cref{sec:actIV-mortgage-pairs}); testing it needs $D{=}1$ pairs
(\Cref{sec:roadmap}). What the box does establish is the miss itself, which is
exactly the deployment risk the dual-label design exists to surface, and exactly what a single
general-safety score reports as clean.

Two numbers are deliberately \emph{not} reported. First, $\text{AP}\cdot G$ rests on only 6 G1
positives in this split and is correspondingly noisy --- a random ranker already scores
$\text{AP}\cdot G\approx 6/146\approx0.04$, so that column must be read against this chance floor, not
against $0$, and the spread from SmolLM2's $0.261$ to Qwen3-4B's $1.000$ is a wide small-sample band, not
four comparable point estimates. Second, we do not report a fixed-threshold
``caught at 5\% FPR'' count for the G0/D1 stratum: for these zero-shot guards the unsafe-probability
distribution is tightly clustered, so a dev-calibrated cutoff sits on a knife-edge --- its G0/D1 catch
count swung by more than 50 rows between library versions of the quantile routine, i.e.\ it is not
reproducible. That instability is itself a finding: \textbf{naive threshold transfer is unreliable} for
these guards, echoing the ranking-recovery-is-not-calibration lesson of Acts I and II
(\Cref{sec:actI-op,sec:actIII}). We therefore report \emph{ranking} (AP) and \emph{invariance}
($\Delta_{\mathrm{context}}$), and leave the operating point to a re-calibration study.

\begin{takeaway}{What ranking sees here, and what it misses}
A general guard, zero-shot, already ranks mortgage violations \emph{moderately} --- ``help me
discriminate'' pattern-matches to unsafe even without a jailbreak, so the domain is not invisible. But
``moderate'' is the whole story: $\text{AP}\cdot D$ tops out at $0.85$, so much of the subtle G0/D1
stratum is unranked --- and against the $0.555$ chance floor set by the $81/146$ $D$-positives, that
band is only $0.12$--$0.30$ above chance. Ranking also says \emph{nothing} about fairness, though on this
split our fairness instrument says little either (\Cref{sec:actIV-mortgage-pairs}). We call $0.85$ ``moderate'' relative to
the $1.0$ of a perfect ranker, not against a known achievable ceiling: we have \emph{no} purpose-built
mortgage-compliance guard on this benchmark to say how high $\text{AP}\cdot D$ could go, so the gap between
$0.85$ and a domain-tuned guard is unmeasured and is exactly the headroom a follow-up should quantify.
Either way the practitioner lesson holds: compliance screening needs a \emph{domain-grounded} ranking metric
\emph{and} a separate invariance gate; a single general-safety score covers neither.
\end{takeaway}

\begin{takeaway}{A measuring stick, not a legal finding}
Four caveats bound everything in this section. \textbf{(1) Not SME-validated}: the labels are
LLM-judge, policy-card-consistent, measured by self-consistency --- no compliance lawyer signed the 24
cards and there is no human Fleiss-$\kappa$. \textbf{(2) Not a full $2\times2$}: the G1/D0 quadrant is
empty, so orthogonality is shown on three quadrants. \textbf{(3) Frozen, not regenerable}: generation is
stochastic and intentionally fixed; only the \emph{evaluation} reproduces. \textbf{(4) Small samples}:
the public-test baselines rest on a single split with 6 G1 positives and 3 protected pairs, so
``fairest'' is a thin claim as stated. Because these guards are \emph{zero-shot} (never tuned on the
benchmark), all \textbf{39} release protected pairs are legitimate evaluation data --- scoring them is a
zero-training recompute that would firm up the invariance ranking, and is listed as a roadmap step
(reported separately from the public-test pairs, so the number stays comparable with a future
fine-tuned arm). The benchmark \emph{surfaces} guard behavior on the G0/D1 stratum and on protected-pair
invariance; it
certifies nothing about any real lender, model, or population. The path to confirmatory use is stated
in \Cref{sec:roadmap}: SME adjudication with per-label Fleiss-$\kappa$, populating G1/D0,
re-decontaminating against the v2 general sources, and adding a domain-fine-tuned guard arm.
\end{takeaway}
